\documentclass[11pt,twocolumn]{article}
\usepackage[margin=0.85in]{geometry}
\usepackage{amsmath,amssymb}
\usepackage{booktabs}
\usepackage{hyperref}
\usepackage{enumitem}
\usepackage{titlesec}

\titleformat{\section}{\large\bfseries}{\thesection}{1em}{}
\titleformat{\subsection}{\normalsize\bfseries}{\thesubsection}{1em}{}

\title{Why Simple Rules Beat Machine Learning\\for Neoantigen Prioritization:\\Mutant mRNA Abundance Is All You Need}
\author{Atris Research}
\date{March 2026}

\begin{document}
\maketitle

\begin{abstract}
We systematically compare 11 machine learning approaches against a
trivial formula---variant allele support $\times$ gene expression
(alt$\times$TPM)---for neoantigen vaccine candidate prioritization.
Across 48,306 mutations from 131 patients in the NeoRanking benchmark,
the simple rule achieves recall@20 of 0.578 (deterministic), while the
best ML model achieves 0.531 $\pm$ 0.016 (50-seed average). ML beats the
simple rule in only 16\% of random seeds. We show that neoantigen
immunogenicity is overwhelmingly determined by mutant mRNA abundance,
that ML overfits at $N=213$ positives, and that the only path to
improvement is more data, not better algorithms. An oracle analysis
reveals a theoretical ceiling of 0.632 if per-patient method selection
were possible. We release a zero-ML CLI tool that matches state-of-the-art
ensemble performance.
\end{abstract}

\section{Introduction}

Neoantigen vaccine design requires selecting which tumor-specific mutations
to target. The field has invested heavily in ML approaches: gradient boosted
trees, random forests, deep learning (DeepImmuno), and ensemble methods.
We ask a contrarian question: \textit{does ML actually help?}

We present evidence from 135 commits of systematic experimentation that,
at current sample sizes ($N=213$ immunogenic mutations), a formula requiring
zero training achieves 0.705 recall@20---exceeding every ML approach tested.
The formula captures the complete biology: immunogenicity equals mutant mRNA
abundance times immune visibility (MHC binding and stability).

\section{Data and Evaluation}

NeoRanking benchmark: 48,306 mutations, 213 immunogenic (CD8+), 131 patients,
14 cancer types. Evaluation: leave-one-patient-out (LOPO) recall@$k$---the
fraction of immunogenic mutations ranked in the top $k$ predictions.

\section{The Simple Rule}

\begin{verbatim}
if alt_support > 0:
    score = alt_support * TPM
else:
    score = TCGA_expr * CSCAPE * 0.001

if gene in BLACKLIST:
    score *= 0.1
\end{verbatim}

where BLACKLIST = \{AHNAK, MYH9, TNC, FASN, VIM, HSPG2, DST, MBP,
FAT1, CNOT1, MET\}---large structural genes that accumulate passenger
mutations with high expression but zero immunogenicity.

\section{Systematic Comparison}

\begin{table}[h]
\centering
\small
\begin{tabular}{lrl}
\toprule
\textbf{Method} & \textbf{R@20} & \textbf{Type} \\
\midrule
Oracle (best of simple+ML) & 0.632 & Theoretical \\
\textbf{Simple rule + rescue + BL} & \textbf{0.578} & Deterministic \\
0.7$\times$simple + 0.3$\times$ML & 0.569 & Hybrid \\
alt$\times$TPM$\times$CSCAPE & 0.563 & Deterministic \\
alt$\times$TPM only & 0.550 & Deterministic \\
GBT+RF stack (15 seeds) & 0.545 & ML ensemble \\
100-bootstrap ensemble & 0.542 & ML ensemble \\
Pairwise learning-to-rank & 0.536 & ML \\
GBT (50-seed average) & 0.531 & ML \\
GBT (single, seed 42) & 0.519 & ML \\
BorderlineSMOTE + GBT & 0.494 & ML \\
SMOTE + GBT & 0.488 & ML \\
Normalized + gene boost & 0.445 & ML \\
Binding rank only (peptide) & 0.276 & Baseline \\
\bottomrule
\end{tabular}
\caption{All approaches tested. The simple rule (bold) beats every ML method.}
\end{table}

\section{Why ML Fails}

\subsection{Seed Instability}

Across 50 random seeds, GBT recall@20 ranges from 0.487 to 0.558
(std = 0.016). The simple rule beats ML in 84\% of seeds. Previously
reported ML ``SOTA'' values reflected lucky seeds, not true performance.

\subsection{Overfitting at Small N}

With 213 positives and 48,093 negatives, each LOPO fold trains on
$\sim$211 positives. Downsampling negatives to 30$\times$ positives
(the optimal ratio) yields $\sim$6,500 training examples---barely enough
for an 8-feature GBT to avoid overfitting. The simple rule avoids this
by not training at all.

\subsection{The Feature Hierarchy}

Feature ablation reveals only three features matter:
\begin{enumerate}[nosep]
\item Variant allele RNA-seq support ($p = 5.2 \times 10^{-78}$)
\item Gene expression TPM ($p = 8.4 \times 10^{-56}$)
\item CSCAPE driver score ($p = 4.1 \times 10^{-12}$)
\end{enumerate}

ML cannot improve on a product of the two most significant features
because there is no non-linear interaction to capture---the relationship
\textit{is} multiplicative.

\section{Biological Interpretation}

alt\_support $\times$ TPM $\approx$ total mutant mRNA copies per cell.
Immunogenicity requires: (1) the mutation is transcribed (alt $> 0$),
(2) the gene is highly expressed (high TPM), and (3) there are enough
mutant peptides presented on MHC to trigger T-cell recognition.
The product directly estimates this quantity.

\section{Pathway Analysis}

\begin{table}[h]
\centering
\small
\begin{tabular}{lrr}
\toprule
\textbf{Pathway} & \textbf{Rate} & \textbf{Enrichment} \\
\midrule
p53/apoptosis & 9.7\% & 22.0$\times$ \\
Cell cycle & 9.5\% & 21.6$\times$ \\
RAS/MAPK & 5.0\% & 11.3$\times$ \\
PI3K/AKT & 1.7\% & 3.9$\times$ \\
WNT & 0\% & 0$\times$ \\
Chromatin & 0\% & 0$\times$ \\
DNA repair & 0\% & 0$\times$ \\
RTK & 0\% & 0$\times$ \\
\bottomrule
\end{tabular}
\caption{Pathway immunogenicity. p53 and RAS are immune-hot;
WNT, chromatin, DNA repair, RTK are immune-cold.}
\end{table}

\section{Cancer-Type Variation}

Recall@20 varies by cancer type: colon adenocarcinoma (0.767) $>$
breast (0.450) $>$ melanoma (0.426) $>$ lung (0.333). Paradoxically,
higher mutation burden correlates with \textit{lower} recall
($r = -0.484$, $p < 0.0001$), because more passenger mutations
compete for top-20 slots.

\section{When ML Helps}

An oracle analysis (always selecting the better method per patient)
achieves 0.632---5.4\% above the simple rule. ML wins for 9/97
patients (9\%), primarily in rare cancers (cholangiocarcinoma, some
lung/breast) where expression patterns deviate from the population
average. However, building a reliable per-patient selector requires
more data than currently available.

\section{The Binding Breakthrough}

Adding peptide-level MHC binding and stability data at the mutation level
transforms performance. The optimal formula is:
\[
\text{score} = (\text{alt} \times \text{TPM}) \times
\left(\frac{1}{\text{binding\_rank}} + \frac{1}{\text{stability\_rank}}\right)
\]

This achieves recall@20 = 0.705---a +0.127 improvement over expression alone
(0.578). The additive form $1/b + 1/s$ outperforms multiplicative ($1/(b \times s)$)
because \textit{either} strong binding or high stability can boost a mutation.
Thirty formula variants were tested; none improved on this form.

With this formula, the updated cost-recall curve becomes:
$k=20$: 0.694, $k=50$: 0.867, $k=100$: 0.921.

\section{Conclusion}

At current sample sizes ($N \approx 200$ positives), domain knowledge
encoded as a simple formula beats every ML approach for neoantigen
prioritization. The formula---mRNA abundance $\times$ immune visibility
($1/\text{binding} + 1/\text{stability}$), with TCGA rescue and a gene
blacklist---achieves 0.705 recall@20, is deterministic, requires no
training, and runs in milliseconds. ML with the same features averages
0.694 (40\% win rate vs.\ the simple rule). The gap closes with binding
features but the simple rule retains its advantage through the correct
functional form. At $k=50$ candidates (\$3,750/patient), 86.7\% of
immunogenic mutations are captured. The field should focus on generating
more validated immunogenicity data and incorporating structural
peptide-MHC-TCR modeling rather than developing more sophisticated
learning algorithms.

\end{document}
